% 統計力学1 統計力学1-2 問題1（001-levels。当時は 統計力学1-1 問題1・shou01-l1-levels） 解いた日 2026-10-27
%% kotae
\toi{1}{◎}{}
$\{n_{x},n_{y},n_{z}\}$

i) $=(1,1,1)$
\[ \frac{E}{\varepsilon_{0}}=3\qquad\text{縮退度：}1 \]
ii) $=(1,1,2)\ (1,2,1)\ (2,1,1)$
\[ \frac{E}{\varepsilon_{0}}=6\qquad\text{縮退度：}3 \]
iii) $=(1,2,2)\ (2,1,2)\ (2,2,1)$
\[ \frac{E}{\varepsilon_{0}}=9\qquad\text{縮退度：}3 \]
iv) $=(1,1,3)\ (1,3,1)\ (3,1,1)$
\[ \frac{E}{\varepsilon_{0}}=11\qquad\text{縮退度：}3 \]
\toi{2}{◎}{}
\[ \varepsilon_{0}=\frac{\pi^{2}\hbar^{2}}{2m_{\mathrm{He}}L^{2}}=\qty{8.26e-38}{\joule} \]
\[ \Delta E=3\varepsilon_{0}=\qty{2.48e-37}{\joule} \]
\toi{3}{◎}{}
\[ \frac{\Delta E}{\kB T}=\num{6.0e-17} \]
準位が変わってもほとんどエネルギーが変化せず，ゆっくりと変わるため。
%% kizuki
